\section{Methodology}\label{sec:methodology}
In this section, we present the end-to-end \textsc{QuaRK} pipeline. In Section~\ref{subsec:QR-embedding}, we introduce our quantum reservoir embedding pipeline. Section~\ref{subsec:K-read} outlines how the reservoir is read out efficiently to form feature vectors and how these are lifted into an RKHS for learning. Section~\ref{subsec:readout-training} describes how the resulting kernel-based readout is classically trained. Finally, Section~\ref{subsec:learning-theoretic-analysis} analyzes, in a PAC-style fashion, the link between design/resource choices and finite-sample performance on weakly dependent data.

\subsection{Quantum reservoir embedding}\label{subsec:QR-embedding}
Although we described the machine learning problem in Section~\ref{subsec:learning-temporal-data} as that concerning learning the mapping $H^\star: \rmX \mapsto Y_0$, in reality we do not have a dataset of time series samples $\rvx \in \gIZ$ to train on. All we have is a single time series, i.e.\ a single realization of the process $\rmX$, which we must leverage to approximate the unknown $H^\star$. To this end, we adopt the windows-based learning approach discussed in~\cite{Gonon_Grigoryeva_Ortega_2020} to form an (artificial) windows dataset $\gD = \{(\rvx_1, y_1), \dots, (\rvx_N, y_N)\}$, where each window is a lookback sequence $\rvx$ of the form $\rvx = (x_{\tau-w+1}, \dots, x_{\tau}) \in (\gI)^w$ along with its label $y = y_\tau$ for some $\tau \in \sZ_-$. We override the definition of the functional $H^T(\cdot)$, which was initially defined to accept semi-infinite inputs, with a definition of $H^T(\rvx)$ for a given window $\rvx$ that starts from an initial state $\rho_{\tau-w}$, evolves the reservoir with respect to inputs $x_t \in \rvx$, and returns the output state $\rho_\tau$ at the end of the evolution.

\subsubsection{Projection}\label{subsubsec:projection} 
Let $n \in \sN$ be the number of qubits chosen by the user and let $\rvx \in \gD$. To make our model parameters independent of the data dimension, which is a key property for allowing scalability with potentially high dimensions $d$, we first project each input $x \in \rvx$ into another representation $z \in \sR^n$ using a Johnson--Lindenstrauss (JL) data projection with a linear projection matrix $\Pi \in \sR^{n \times d}$ so that $z = \Pi x$. We use a random matrix projection $\Pi$ distributed according to the Gaussian JL distribution~\cite{Dasgupta_2002}
\begin{equation}\label{eq:JL-matrix}
    \Pi = \pth{\Pi_{ij}}_{i,j} \in \sR^{n \times d}, \qquad \Pi_{ij} \overset{\mathrm{i.i.d.}}{\sim} \gN\pth{0, \tfrac{1}{n}}.
\end{equation}
Such a projection is deliberately chosen to yield low-distortion embeddings, so that the (Euclidean) distances between all points in the dataset---i.e.\ the set $\gP := \bigcup_{i=1}^N \rvx_i \subset \sR^d$ of cardinality at most $wN$---are nearly preserved. Importantly, one can show that if the number of qubits satisfies $n = \Omega\big(\varepsilon_{\rm pr}^{-2} \log\pth{\tfrac{w N}{\delta}}\big)$ for some error tolerance and failure probability $\varepsilon_{\rm pr}, \delta \in (0,1)$, then for all $u,v \in \gP$, we have
\begin{equation}
    (1-\varepsilon_{\rm pr}) \normeuc{u - v}^2 \leq \normeuc{\Pi u - \Pi v}^2 \leq (1+\varepsilon_{\rm pr}) \normeuc{u - v}^2,
\end{equation}
with probability at least $1-\delta$ by the standard JL lemma (see Theorem 2.1 in~\cite{woodruff2014sketching}). This observation is later leveraged to show effective learning with our reservoir computer (see Section~\ref{subsec:effective-learning}).

\subsubsection{Quantum reservoir architecture}
For the QR architecture, we adopt a design that adheres to the family of quantum reservoirs called contracted-encoding quantum channels (CEQC) as introduced in~\cite{Martínez-Peña_Ortega_2025}. Such a design is interesting for numerous reasons, the main one being its automatic satisfaction of the FMP and ESP. These channels are defined as CPTP maps $T: \gSH \times \gI \to \gSH$ that lead to state-space transformations of the form
\begin{equation}\label{eq:CEQC}
    \rho_t = T(\rho_{t-1}, x_t) := \gE \circ \gJ(\rho_{t-1}, x_t),
\end{equation}
where $\gE: \gSH \to \gSH$ is a strictly contractive map, and $\gJ: \gSH \times \gI \to \gSH$ is a CPTP map that encodes the input information. It is worth mentioning that quantum reservoirs defined by CPTP maps of this form encompass numerous architectures in the QRC literature, such as quantum circuits with noise~\cite{Suzuki_2022, Kubota_2023}, reset-rate channels~\cite{Chen_Nurdin_Yamamoto_2020, Molteni_2023}, mid-circuit measurements~\cite{Yasuda_Suzuki_Kubota_Nakajima_Gao_Zhang_Shimono_Nurdin_Yamamoto_2023, Hu_2024}, and compositions of a dissipative channel with an amplitude encoding map~\cite{Mujal_2023}. The following paragraphs present our design for both channels $\gJ$ and $\gE$.

\paragraph{Unitary evolution.}
After projecting an input $x$ into $z = \Pi x$, we inject each component $z_j$ of $z$ into the circuit using $R_y$ gates parametrized by angles $\theta(z_j) = \pi \cdot \tanh(z_j) \in [-\pi, \pi]$, followed by a global fixed unitary $W$, which realizes the evolution $\gJ(\rho,x) = V(x) \rho V(x)^\dagger$, where
\begin{equation}\label{eq:unitary-design}
    V(x) := W \bigotimes_{j=1}^n R_y\big(\theta(z_j)\big)
\end{equation}
as illustrated in Figure~\ref{fig:unitary-evolution}. Furthermore, given our quantum hardware topology $(V, E)$, where the vertices are our $n$ qubits and $E$ contains available edges between qubits of the form $(i,j)$, we realize the unitary $W$ following the Ising-like, hardware-friendly design
\begin{equation}\label{eq:ising-unitary}
    W = \prod_{j \in V} R_x(\vartheta_{x}^j) \prod_{j \in V} R_z(\vartheta_{z}^j) \prod_{(i,j) \in E} R_{zz}(\vartheta_{zz}^{ij}),
\end{equation}
where all parameters $\vartheta_{\bullet}$ are sampled i.i.d.\ uniformly from $(-\pi, \pi)$.

\begin{figure}[t]
\centering
\scalebox{0.8}{
% \begin{quantikz}[column sep=0.45cm]
\begin{quantikz}
\lstick{$q_1$} &
\qw \gategroup[wires=3,steps=10,
  style={dashed, rounded corners, fill=gray!10, inner xsep=4pt, inner ysep=17pt},
  background,
  label style={label position=below, anchor=north, yshift=-0.25cm}
]{Unitary $V(x)$} & \gate{R_y(\theta(z_1))} &
  \ctrl[wire style={"ZZ(\vartheta_{zz}^{12})"}]{1}
    \gategroup[wires=3,steps=7,
      style={dashed,rounded corners,fill=blue!20,inner xsep=4pt,inner ysep=3pt},
      background,
      label style={label position=below,anchor=north,yshift=-0.2cm}
    ]{\small Ising-like $W$}
  & & &
  \ctrl[wire style={"ZZ(\vartheta_{zz}^{31})"{right,yshift=-7pt}}]{2} & &
  \gate{R_z(\vartheta^1_z)} &
  \gate{R_x(\vartheta^1_x)} & \qw & \slice[style={red,dashed},label style={red}]{\(\rho\)} &
  \qw \\
\lstick{$q_2$} & \qw &
  \gate{R_y(\theta(z_2))} &
  \control{} &
  \ctrl[wire style={"ZZ(\vartheta_{zz}^{23})"{right,yshift=-5pt}}]{1} & & & &
  \gate{R_z(\vartheta^2_z)} &
  \gate{R_x(\vartheta^2_x)} & \qw & 
  \qw & \\
\lstick{$q_3$} & \qw &
  \gate{R_y(\theta(z_3))} & &
  \control{} & &
  \control{} & &
  \gate{R_z(\vartheta^3_z)} &
  \gate{R_x(\vartheta^3_x)} & \qw &
  \qw &
\end{quantikz}
}
\caption{\small Circuit for the unitary evolution block $V(x)$ on three qubits with ring connectivity.
First, classical features are encoded via single-qubit angle encoding $R_y(\theta(z_j))$ on each $q_j$.
Then the Ising-like unitary $W$ (~\cref{eq:ising-unitary}) is shown as ZZ couplings on the edges $(1,2)$, $(2,3)$, and $(3,1)$,
followed by local rotations $R_z(\vartheta_z)$ and $R_x(\vartheta_x)$.}
\Description{Schematic quantum circuit on three qubits q1--q3 showing an input-dependent unitary block. Each qubit first undergoes a single-qubit Y-rotation Ry(theta(z_j)) to encode classical features. A highlighted inner region then applies three pairwise ZZ couplings arranged in a ring (between qubit pairs 1--2, 2--3, and 3--1). The block ends with local single-qubit rotations Rz and Rx on each qubit. A dashed outer box indicates the overall unitary V(x), and the highlighted region marks the contractive channel component within it.}

\label{fig:unitary-evolution}
\end{figure}


\paragraph{Contractive quantum channel $\gE$.} Let $\lambda \in (0,1)$ be a contraction factor. We design the channel $\gE$ as the reset-rate channel
\begin{equation}\label{eq:contractive-channel}
    \gE_\lambda: \gSH \longto \gSH, \quad \rho \longmapsto \lambda \rho + (1 - \lambda) \pdensity.
\end{equation}
As it is easy to check (with $\gE_\lambda(\rho')$ in the first line below) that
\begin{equation}
\begin{aligned}
    &\normHS{\gE_\lambda(\rho) - \gE_\lambda(\rho)} \leq \lambda \normHS{\rho - \rho'} \\
    \text{and} \;\; &\norm{\gJ(\rho, x) - \gJ(\rho', x)} \leq \norm{\rho - \rho'},
\end{aligned}
\end{equation}
we have $\normHS{T(\rho, x) - T(\rho', x)} \leq \lambda \normHS{\rho - \rho'}$. Hence, we adhere to the condition in the last paragraph of Section~\ref{subsec:qrc-background}, thereby making our reservoir computer satisfy both the ESP and FMP\footnote{Notice also that our map $T = \gE_\lambda \circ \gJ$ is by definition non-unital, i.e., $T(I, x) \neq I$, hence we ensure a non-pathological, non-trivial filter as discussed in Theorem 1 of~\cite{Mart_nez_Pe_a_2023}.}. Figure~\ref{fig:contractive-channel} illustrates how we realize such a reset-rate channel on our quantum circuit, where, controlled on a coin qubit $c$ rotated by $R_y(\thetal)$ with $\thetal = 2 \arcsin \sqrt{1-\lambda}$, we swap reservoir qubits with ancilla ones all initialized to $\ket{+}$. The convex formula in~\cref{eq:contractive-channel} is thus realized, as the coin $c$ becomes superposed with probability $1 - \lambda$ of being in state $\ket{1}$. In what follows, we combine parameters $\vartheta_{\bullet}$ and $\lambda$ into a single vector $\vartheta$.

\paragraph{Quantum embedding of a window.}\label{par:quantum-embedding} Let $\rvx = (x_{\tau-w+1}, \ldots, x_{\tau})$ be a window of our time series. A consequence of using the reset-rate channel $\gE_\lambda$ is that, provided $w$ is large enough, the state at which our quantum reservoir is initialized at time $\tau-w$ matters little to the output state $\rho_\tau$ (see the third paragraph of Section~\ref{subsec:qrc-background}). Hence, by convention, we fix $\rho_{\tau-w} = \pdensity$ as the initial state of the reservoir at time $\tau-w$, and we process the inputs $x_t$ step by step until reaching the output state $\rho_\tau$, realizing the following reservoir evolution:
\begin{equation}
    \rho_{\tau-w} \xlongrightarrow{T(\cdot, \, x_{\tau - w + 1})} \rho_{\tau-w+1} \xlongrightarrow{T(\cdot, \, x_{\tau - w + 2})} \cdots \xlongrightarrow{T(\cdot, \, x_{\tau})} \rho_{\tau}.
\end{equation}
We refer to such a window embedding as $H^T(\rvx) \in \gSH$, which maps the window $\rvx$ to $\rho_\tau$.

\begin{figure}[t]
\centering
\scalebox{1}{
\begin{quantikz}[row sep=0.3cm]
\lstick{$q_1$} &
\qw
\gategroup[wires=7,steps=5,
  style={dashed, rounded corners, fill=gray!10, inner xsep=4pt, inner ysep=3pt},
  background,
  label style={label position=below, anchor=north, yshift=-0.25cm}
]{Contractive channel $E_\lambda$}
& \qw & \swap{3} & \qw      & \qw      & \qw \\[0.2cm]
\lstick{$q_2$} & \qw & \qw & \qw      & \swap{3} & \qw      & \qw \\[0.2cm]
\lstick{$q_3$} & \qw & \qw & \qw      & \qw      & \swap{3} & \qw \\[0.2cm]
\lstick{$a_1$} & \gate{\ket{0}} & \gate{H} & \swap{-3} & \qw      & \qw      & \qw \\
\lstick{$a_2$} & \gate{\ket{0}} & \gate{H} & \qw      & \swap{-3} & \qw      & \qw \\
\lstick{$a_3$} & \gate{\ket{0}} & \gate{H} & \qw      & \qw      & \swap{-3} & \qw \\
\lstick{$c$}   & \gate{\ket{0}} & \gate{R_y(\theta_\lambda)} & \ctrl{-6} & \ctrl{-5} & \ctrl{-4} & \qw
\end{quantikz}
}

\caption{\small SWAP-dilation realization of the contractive channel $E_\lambda$ acting on a 3-qubit reservoir state $\rho$ (top wires).
The ancilla register (middle wires) is reset to $\ket{0}$ and prepared in $\ket{+}^{\otimes 3}$ via Hadamards.
A coin qubit (bottom wire) is reset and rotated by $R_y(\theta_\lambda)$, where $\theta_\lambda = 2\arcsin\!\sqrt{1-\lambda}$ so that
$\Pr[c=1]=1-\lambda$ and $\Pr[c=0]=\lambda$.
Conditioned on $c=1$, controlled-SWAPs exchange each reservoir qubit $q_i$ with its corresponding ancilla qubit $a_i$.}
\Description{Circuit diagram implementing a contractive channel via a SWAP-dilation. The top three wires are reservoir qubits q1--q3. Three ancilla wires a1--a3 are initialized to |0> and prepared with Hadamard gates. A bottom “coin” qubit c is initialized to |0> and rotated by Ry(theta_lambda). Controlled-SWAP gates, controlled by the coin qubit, swap each reservoir qubit qi with its corresponding ancilla ai when the control is active. A dashed box encloses the full channel implementation.}
\label{fig:contractive-channel}
\end{figure}

\subsubsection{Increased expressivity via spatial multiplexing}
As introduced in~\cite{Nakajima_Fujii_Negoro_Mitarai_Kitagawa_2019} and revised in~\cite{Chen_Nurdin_2019,Chen_Nurdin_Yamamoto_2020}, we adopt the idea of spatial multiplexing to increase model expressivity. The idea is that, instead of using a single reservoir computer, we can use a set of $R \geq 1$ sub-reservoirs with the same architecture but different parameters. These sub-reservoirs are evolved independently (e.g., by running different quantum circuits with different parameters $\vartheta$), so that for a single window $\rvx$ we obtain a set of $R$ different quantum states, thereby increasing the feature vector size and expressivity. We mathematically represent the resulting spatially multiplexed (SM) window embedding using the compact equivalent form 
\begin{equation}
    H^T: \rvx = (x_{\tau-w+1}, \dots, x_{\tau}) \longmapsto \bigotimes_{r=1}^R H^{T_r}(\rvx) := \rho_{\tau}^{(r)} \in \gSH^R,
\end{equation}
with $T_r$ being the evolution map of sub-reservoir $r \in [R]$.
In addition, the parameters $\vartheta_r$ are all sampled independently. Working with different contraction values $\lambda_r$ is beneficial, as it yields different sub-functionals $H^{T_r}$ each with a different fading-memory structure, thereby enriching model expressivity with respect to modeled memories\footnote{It has been shown that the family of QRCs is universal in that of the functionals with the FMP~\cite{Nakajima_Fujii_Negoro_Mitarai_Kitagawa_2019}.}.


\subsection{Reading out the reservoir}\label{subsec:K-read}
In this section, we describe the readout module. We first present an efficient classical shadows measurement scheme to form feature vectors from the reservoir state, then explain how a kernel-based RKHS readout is trained on these features.

\subsubsection{Efficient measurement scheme}\label{subsubsec:efficient-msmt-scheme}
At the end of the reservoir, we read out each sub-reservoir state $\rho_\tau$ by measuring a set of $k$-local observables denoted by $\gO$---in our experiments in Section~\ref{sec:experiments}, we set $\gO$ to be the set of $2$-local Pauli observables\footnote{A Pauli observable $P =\sigma_1 \otimes \cdots \otimes \sigma_n$, with $\sigma_i \in \{I, X, Y, Z\},$ is called $k$-local (or weight $\leq k$) if the index set $\{i \in [n]: \sigma_i \neq I\}$ has cardinality $\leq k$.}. Such a choice is deliberately made so that measuring the sub-reservoirs can be efficiently realized. Indeed, we make use of the classical-shadow estimator~\cite{huang-classical-shadows-2020}, which states that one can simultaneously and efficiently estimate all moments $[\braket{O}_{\rho_\tau}: O \in \gO]$ up to error tolerance $\varepsilon_{\rm cs} \in (0,1)$ using only a number of circuit runs (classical snapshots) scaling as $O\big(\pth{\tfrac{3}{2}}^k \varepsilon_{\rm cs}^{-2} \log(|\gO|)\big)$ by employing random single-qubit Pauli basis measurements~\cite{huang-classical-shadows-2020}.

The moments are estimated using the median-of-means algorithm, as introduced in the seminal paper on classical shadows~\cite{huang-classical-shadows-2020}, which is more resistant to statistical fluctuations. In the end, our measurement scheme for the sub-reservoirs $m: \gSH^R \to \sR^{R |\gO|}$ produces the feature vector
\begin{equation}\label{eq:feature-vector}
    \Phi(\rvx) = m \pth{H^T(\rvx)} := \cro{\braket{O}_{\rho_\tau^{(r)}}: O \in \gO, \;  r \in [R]}^\top \in \sR^{R |\gO|},
\end{equation}
where $H^T(\rvx) = \bigotimes_r \rho_\tau^{(r)}$.
Subsequently, feature vectors of this kind form the inputs to the classical Mat\'ern kernel we use in the following section.

\subsubsection{RKHS readout}
Let $\varphi$ be a (unit-variance) Mat\'ern covariance with smoothness and lengthscale parameters $\nu > 1$ and $\xi > 0$, respectively, defined as
\begin{equation}\label{eq:matern-profile}
    \varphi(s) := \frac{2^{1-\nu}}{\Gamma(\nu)} \pth{\sqrt{2 \nu}\; \frac{s}{\xi}}^\nu K_\nu\pth{\sqrt{2\nu} \; \frac{s}{\xi}},
\end{equation}
where $\Gamma$ is the gamma function, and $K_\nu$ is the modified Bessel function of the second kind~\cite{genton2001classes}. We then define our kernel function $\kappa$ on $\gSH^R$ as
\begin{equation}\label{eq:matern-kernel}
    \kappa: 
    \begin{array}{rcl}
    \gSH^R \times \gSH^R &\longto& \sR\\
    (\rho, \rho') &\longmapsto& 
        \varphi \Big(\normeuc{m(\rho)-m(\rho')} \Big).
    \end{array}
\end{equation}
Finally, the classical readout functions $h$ we use are drawn from the reproducing kernel Hilbert space (RKHS) $\gHkp$ of $\kappa$ defined as the closure (inclusion of limit points) of the set of functions~\citep{Hofmann_Schölkopf_Smola_2008}
\begin{equation}
    \gHkp^{(0)} := \Set{h = \sum_{i=1}^{s_h} \alpha_i \kappa(\rho_i, \cdot) \;:\;  \alpha_i \in \sR, \; \rho_i \in \gSH, \; s_h \in \sN},
\end{equation}
i.e., $\gHkp := \overline{\gHkp^{(0)}}$. We equip such an RKHS space with the norm defined for each $h \in \gHkp$ as $\normkp{h} := \sqrt{\Braket{h, h}_\kappa}$, with $\Braket{\cdot, \cdot}_\kappa$ being the RKHS's inner-product~\cite{Hofmann_Schölkopf_Smola_2008}.

This readout scheme adheres to the principle of projected quantum kernels~\cite{Huang_Broughton_Mohseni_Babbush_Boixo_Neven_McClean_2021}, consisting of first measuring and then applying a kernel wrapper to the resulting feature vectors. Hence, our approach is a legitimate quantum kernel method applied to time series data. Additionally, we choose a Mat\'ern profile because its RKHS exhibits a polynomial algebraic structure, which is required to preserve the polynomial-algebra structure induced by the spatial multiplexing design used to increase model expressivity~\cite{Monzani_Prati_2025}.

\subsection{Readout training with empirical risk minimization}\label{subsec:readout-training}
Given a fixed quantum channel $T$, i.e.\ after sampling parameters $\vartheta$, and fixing the number of qubits $n$ and the number of sub-reservoirs $R$, we now show how the kernel-based readout $h$ is trained. 
Let $\varphi$ be a Mat\'ern profile, as described in~\cref{eq:matern-profile}, parameterized by hyperparameters $\nu$ and $\xi$. Suppose we have a (training) dataset of windows $\gD = \{(\rvx_1, y_1), \dots, (\rvx_N, y_N) \}$ that contains window samples $\rvx_i$ along with their labels $y_i \in \sR$, obtained from some unknown functional $H^\star$. We learn the best readout $h \in \gHkp$ following an empirical risk minimization scheme, where we minimize the empirical risk $\hat{R}_N(H^T_h)$ with respect to the dataset $\gD$ given by
\begin{equation}\label{eq:empirical-risk-proxy}
    \hat{R}_N(H^T_h) := \frac{1}{N} \sum_{i=1}^N \ell\pth{H^T_h(\rvx_i), y_i},
\end{equation}
with $\ell$ being the squared loss function\footnote{Other Lipschitz losses may be used as discussed in Section~\ref{subsec:learning-temporal-data}.} $\ell(\hat{y}, y) := (\hat{y} - y)^2$.

The bias-variance tradeoff of the model is controlled by restricting the RKHS readout. Concretely, $h$ is learned by minimizing the empirical risk while enforcing a norm budget $\normkp{h} \leq \Lambda$. Such a constrained problem admits an equivalent Tikhonov-regularized formulation~\cite{scholkopf2001learning} where there exists a $\lambda_{\rm \reg} \geq 0$ (a Lagrange multiplier for the constraint $\normkp{h} \le \Lambda$) for which the solution of the norm-constrained empirical risk minimization satisfies
\begin{equation}
    \argmin_{h\in \gHkp(\Lambda)} \hat{R}_N(H^T_h) = \argmin_{h \in \gHkp} \hat{R}_N(H^T_h) + \lambda_{\rm \reg} \normkp{h}^2,
\end{equation}
where $\gHkp(\Lambda)$ is the RKHS ball of radius $\Lambda$ defined as 
\begin{equation}
    \gHkp(\Lambda) := \{h \in \gHkp : \normkp{h} \leq \Lambda\}
\end{equation}
By the representer theorem~\cite{Hofmann_Schölkopf_Smola_2008}, the optimal readout admits the form $h(\cdot) = \sum_{i=1}^N \alpha_i \kappa\pth{H^T(\rvx_i), H^T(\cdot)}.$ Training thus reduces to kernel ridge regression, where we solve a finite-dimensional linear system in $\valpha$. The closed-form solution is given by
\begin{equation}
    \valpha = \pth{K + N \lambda I}^{-1} \vy,
\end{equation}
where $K$ is the Gram matrix with elements $K_{ij} := \kappa\big(H^T(\rvx_i), H^T(\rvx_j)\big)$ for the different windows $\rvx_i, \rvx_j \in \gD$, and $\vy := (y_1, \ldots, y_N)^\top$ is the label vector.
